% answers_short.tex —— 2015 年·天津卷·理科 · 简短答案速查
% 由 tools/make_short_answers.py 自 answers.tex 的 \gksolhead/\gkanswer 提取，
% 并与 2015/2015_数学真题_天津理_答案要点.md 逐题比对后生成，勿手改（改源 MD 后重跑）。

\begin{gknotice}
\textbf{使用说明}：本节只给最终答案便于快速核对，完整推导见解析版。答案为编写者独立推导，\textbf{非官方标准答案}。
\end{gknotice}

\gkpart{选择题}
\begin{center}\begin{tabular}{ccccccccc}
\toprule
题号 & 1 & 2 & 3 & 4 & 5 & 6 & 7 & 8 \\
\midrule
答案 & A & C & B & A & A & D & C & D \\
\bottomrule
\end{tabular}\end{center}

\gkpart{填空题}
\begin{itemize}[leftmargin=2.6em]
  \item[\textbf{9.}] \(-2\)
  \item[\textbf{10.}] \(\dfrac{8\pi}{3}\ \mathrm{m}^3\)
  \item[\textbf{11.}] \(\dfrac{1}{6}\)
  \item[\textbf{12.}] \(\dfrac{15}{16}\)
  \item[\textbf{13.}] \(8\)
  \item[\textbf{14.}] \(\dfrac{29}{18}\)
\end{itemize}

\gkpart{解答题}
\begin{itemize}[leftmargin=2.6em]
  \item[\textbf{15.}] (1) \(\pi\)；(2) 最大值 \(\dfrac{\sqrt{3}}{4}\)，最小值 \(-\dfrac{1}{2}\)
  \item[\textbf{16.}] (1) \(\dfrac{6}{35}\)；(2) \(P(X=1)=\dfrac{1}{14}\)、\(P(X=2)=\dfrac{3}{7}\)、\(P(X=3)=\dfrac{3}{7}\)、\(P(X=4)=\dfrac{1}{14}\)，数学期望 \(E(X)=\dfrac{5}{2}\)
  \item[\textbf{17.}] (1) 成立（建系后 \(\overrightarrow{MN}=\left(-\dfrac52,0,0\right)\) 与底面法向量垂直、\(z\) 分量为 0，即 \(MN\parallel\) 平面；又 \(M\notin\) 平面 \(ABCD\)，故 \(MN\parallel\) 平面 \(ABCD\)）；(2) \(\dfrac{3\sqrt{10}}{10}\)；(3) \(\sqrt{7}-2\)
  \item[\textbf{18.}] (1) \(q=2\)，\(a_n=2^{\frac{n-1}{2}}\)（\(n\) 为奇数时），\(a_n=2^{\frac{n}{2}}\)（\(n\) 为偶数时）；(2) \(S_n=4-\dfrac{n+2}{2^{n-1}}\)
  \item[\textbf{19.}] (1) \(\dfrac{\sqrt{3}}{3}\)；(2) \(\dfrac{x^2}{3}+\dfrac{y^2}{2}=1\)；(3) \(\left(-\infty,-\dfrac{2\sqrt{3}}{3}\right)\cup\left(\dfrac{\sqrt{2}}{3},\dfrac{2\sqrt{3}}{3}\right)\)
  \item[\textbf{20.}] (1) 当 \(n\) 为偶数时，\(f(x)\) 在 \((-\infty,1)\) 上单调递增、在 \((1,+\infty)\) 上单调递减；当 \(n\) 为奇数时，\(f(x)\) 在 \((-\infty,-1)\) 上单调递减、在 \((-1,1)\) 上单调递增、在 \((1,+\infty)\) 上单调递减；(2) 成立（记 \(t=n^{\frac{1}{n-1}}\)，\(P(t,0)\)；切线 \(g(x)=n(1-n)(x-t)\)，令 \(h=g-f=x^n-n^2x+n(n-1)t\)，则 \(h'=n\left(x^{n-1}-t^{n-1}\right)\)，\(h\) 在 \(x=t\) 取最小值 \(h(t)=0\)，故对任意正实数 \(x\) 有 \(h\geqslant0\)，等号仅当 \(x=t\)）；(3) 成立（记 \(t=n^{\frac{1}{n-1}}\leqslant2\)，等号仅 \(n=2\)；由切线不等式得严格界 \(t-x_2\geqslant\dfrac{a}{n(n-1)}\)，又 \(x_1>\dfrac{a}{n}\)（严格），相减得 \(x_2-x_1<t-\dfrac{a}{n-1}\leqslant2-\dfrac{a}{n-1}=\dfrac{a}{1-n}+2\)）
\end{itemize}
